\documentclass[10pt,a4paper]{article}
\usepackage[utf8]{inputenc}
\usepackage[T1]{fontenc}
\usepackage[french]{babel}
\usepackage[a4paper,margin=1.8cm]{geometry}
\usepackage{amsmath,booktabs,tabularx,array,float,hyperref}
\hypersetup{colorlinks=true,linkcolor=blue,urlcolor=blue}
\newcolumntype{Y}{>{\raggedright\arraybackslash}X}

\title{Rapport de reproduction des baselines\\Phase 2/10}
\author{Stéganographie adaptative par parcours de graphe}
\date{13 juin 2026}

\begin{document}
\maketitle

\begin{abstract}
Cette phase établit les baselines nécessaires à l'évaluation future :
BIND, AdaBIND, BYNIS, marche uniforme, DeepWalk et node2vec. Le code officiel
de Lee et al. a été figé au commit
\texttt{b7499454\ldots539b02fa}. La reproduction de BIND
produit exactement le même ordre d'arêtes que l'oracle HMG. Dix tests
automatisés contrôlent le décodage, la capacité, le déterminisme et la
préservation ou la modification de la topologie. Les résultats obtenus sur
le réseau réel \emph{terrorists-911} confirment la capacité de BIND, mais
mettent également en évidence sa fragilité à de faibles réordonnancements.
\end{abstract}

\section{Périmètre et protocole}

Le jeu de contrôle contient 62 sommets et 152 arêtes dirigées. Son empreinte
SHA-256 commence par \texttt{0897b4df}\allowbreak\texttt{9fcddbb3} et se
termine par \texttt{54d55315}\allowbreak\texttt{bfeb03a4}; sa valeur complète
est archivée dans la sortie JSON.
Il s'agit d'un réseau réel distribué avec HMG. Il sert exclusivement à la
reproduction des baselines : il ne remplace pas les trajectoires et sessions
observées exigées pour la validation de notre méthode.

La graine globale est 1234. BIND conserve le multiensemble des arêtes et
encode les couples de bits par les parités des degrés de leurs extrémités.
AdaBIND ajoute des arêtes lorsque certaines classes sont insuffisantes.
BYNIS construit un graphe synthétique dont la somme des identifiants modulo
256 code un octet. Les sorties sont produites automatiquement aux formats
JSON et CSV par \texttt{code/scripts/run\_phase2\_reproduction.py}.

\section{Conformité à l'implémentation officielle}

Le format de longueur, les quatre classes \(00,01,10,11\), la sélection des
arêtes et la permutation MT19937 sont reproduits. Un test oracle exécute HMG
et notre implémentation sur les mêmes 152 arêtes, le même message de 16 octets
et la même clé. Les deux listes ordonnées sont strictement égales.

Les effectifs observés sont \(51,39,37,25\). La borne équilibrée de Lee vaut
donc
\[
  C_{\mathrm{BIND}}=8\min(51,39,37,25)=200\ \text{bits}.
\]
Cette borne est conservatrice : l'encodabilité réelle dépend de la fréquence
des quatre couples de bits dans le message, en incluant l'en-tête.

\section{Résultats passifs}

\begin{table}[H]
\centering
\caption{Résultats de reproduction sur \emph{terrorists-911}.}
\begin{tabular}{lrrrrl}
\toprule
Méthode & Message & Arêtes finales & Arêtes ajoutées & BPE & Décodage\\
\midrule
BIND & 8 o  & 152 & 0  & 0,421 & exact\\
BIND & 16 o & 152 & 0  & 0,842 & exact\\
BIND & 24 o & 152 & 0  & 1,263 & exact\\
AdaBIND & 8 o & 159 & 7 & 0,403 & exact\\
BYNIS & 18 o & 20 & -- & 7,200 & exact\\
BYNIS + 50 leurres & 18 o & 70 & 50 & 2,057 & exact\\
\bottomrule
\end{tabular}
\end{table}

Le cas AdaBIND utilise huit octets \texttt{ff}, volontairement défavorables à
BIND. Sept arêtes sont ajoutées en sept itérations. Ce succès a un coût
conceptuel majeur : la topologie n'est plus préservée. BYNIS obtient un débit
élevé parce qu'il synthétise le support; il doit rester dans une catégorie
séparée des méthodes utilisant un graphe réel inchangé.

\section{Robustesse au réordonnancement}

Chaque intensité comporte 100 attaques indépendantes sur le message BIND de
16 octets. Une attaque échange des positions adjacentes après émission.

\begin{table}[H]
\centering
\caption{Effet des échanges adjacents sur BIND.}
\begin{tabular}{rrrrr}
\toprule
Échanges & Positions touchées & BER moyenne & Récupération exacte & Échec décodage\\
\midrule
1  & 1,32\%  & 0,116 & 42\% & 5\%\\
2  & 2,63\%  & 0,161 & 24\% & 9\%\\
4  & 5,26\%  & 0,320 & 3\%  & 20\%\\
8  & 10,53\% & 0,506 & 0\%  & 30\%\\
16 & 21,05\% & 0,846 & 0\%  & 57\%\\
\bottomrule
\end{tabular}
\end{table}

L'avis est tranché : BIND est une baseline de capacité utile, mais ce n'est
pas une base suffisamment robuste pour notre proposition. Même une
perturbation minime peut altérer l'en-tête et provoquer une désynchronisation
globale. Notre architecture devra séparer synchronisation, détection d'erreur
et correction, puis les évaluer sous un modèle de canal explicite.

\section{Parcours de référence}

Sur le graphe non dirigé associé, des marches de longueur 64 ont été générées.
La marche uniforme visite 31 sommets distincts et présente un taux de retour
immédiat de 0,210. Avec \(p=0,5\) et \(q=2\), node2vec visite 21 sommets et son
taux de retour atteint 0,306. Le corpus DeepWalk contient 248 marches, toutes
de longueur 64, avec un ratio moyen de sommets distincts de 0,423.

Ces chiffres ne mesurent pas encore la furtivité. Ils vérifient les mécanismes
et montrent que les paramètres de node2vec modifient effectivement la
dynamique locale. La phase de données réelles devra comparer ces parcours à
des trajectoires observées, et non les déclarer naturels par construction.

\section{Corrections et limites}

AdaBIND recalcule les classes après chaque ajout, puisque la parité des deux
extrémités reclassifie aussi leurs arêtes incidentes. Le meilleur candidat est
réinitialisé à chaque itération pour éviter un blocage possible du code
officiel sur une ancienne arête. Ces écarts sont documentés et testés.

La phase ne reproduit pas les campagnes OGB à grande échelle ni les coûts
mémoire publiés. Elle valide les contrats algorithmiques et la petite
expérience réelle distribuée par les auteurs. Les résultats synthétiques de
BYNIS restent des résultats de baseline et ne seront pas utilisés comme preuve
principale dans l'article ou le mémoire.

\section{Décision de passage}

La phase 2 est terminée pour son périmètre :
\begin{itemize}
  \item dix tests automatisés réussissent;
  \item BIND est identique à l'oracle officiel sur le cas contrôlé;
  \item BIND, AdaBIND et BYNIS récupèrent exactement leurs messages;
  \item les effets topologiques sont explicitement séparés;
  \item les trois parcours de référence sont déterministes et testés;
  \item la campagne, la configuration et les résultats sont figés.
\end{itemize}

La phase suivante doit construire le pipeline causal des datasets réels,
leurs fiches de provenance et leurs statistiques temporelles.

\end{document}
